\section{Evaluation}


\subsection{Developer Productivity}

We measured development time for implementing common commerce features with Astle.js versus vanilla Redux:

\begin{table}[h]
\centering
\caption{Development time comparison (hours)}
\label{tab:devtime}
\begin{tabular}{lrr}
\hline
Feature & Redux & Astle.js \\
\hline
Cart management & 16 & 4 \\
Checkout flow & 24 & 8 \\
Real-time inventory & 12 & 2 \\
Offline support & 20 & 4 \\
\hline
Total & 72 & 18 \\
\hline
\end{tabular}
\end{table}

Astle.js reduces implementation time by 75\% through domain-specific abstractions.

\subsection{Bundle Size}

\begin{table}[h]
\centering
\caption{Bundle size comparison (KB, gzipped)}
\label{tab:bundle}
\begin{tabular}{lr}
\hline
Library & Size \\
\hline
Astle.js core & 8.2 \\
Astle.js + React & 11.4 \\
Redux + middleware & 6.8 \\
Custom commerce logic & 15.2 \\
\hline
\end{tabular}
\end{table}

Astle.js bundle (11.4 KB) is smaller than equivalent Redux setup (22 KB) while providing more functionality.

\subsection{Runtime Performance}

We benchmarked state update performance:

\begin{table}[h]
\centering
\caption{Operations per second (higher is better)}
\label{tab:perf}
\begin{tabular}{lrr}
\hline
Operation & Redux & Astle.js \\
\hline
Simple dispatch & 125,000 & 98,000 \\
Cart add item & 45,000 & 52,000 \\
Full checkout & 8,200 & 9,100 \\
\hline
\end{tabular}
\end{table}

Astle.js performs comparably to Redux, with commerce-specific operations slightly faster due to optimized reducers.

\subsection{Case Study: Production Deployment}

We deployed Astle.js to a fashion retailer's storefront (50,000 daily active users). Key metrics:

\begin{itemize}
\item Cart abandonment reduced 12\% (optimistic updates reduced perceived latency).
\item Mobile conversion increased 8\% (offline support enabled subway shopping).
\item Support tickets reduced 23\% (time-travel debugging simplified issue reproduction).
\end{itemize}

